\chapter{集合和逻辑用语}

%% ---- 知识框架 ---- %%
\begin{knowledgeframe}
\begin{itemize}
	\item \textbf{集合的基本概念}：元素与集合的关系（$\in, \notin$）、集合的表示法（列举法、描述法）、常用数集符号（$\N, \Z, \Q, \R$）。
	\item \textbf{集合间的关系}：子集（$\subseteq$）、真子集（$\subsetneq$）、相等（$=$）、空集（$\emptyset$）。
	\item \textbf{集合的基本运算}：交集（$\cap$）、并集（$\cup$）、补集（$\complement_U A$）、差集（$A \setminus B$）。
	\item \textbf{运算律}：德摩根定律——$\complement_U(A \cup B) = (\complement_U A) \cap (\complement_U B)$，\ $\complement_U(A \cap B) = (\complement_U A) \cup (\complement_U B)$。
	\item \textbf{容斥原理}：$|A \cup B| = |A| + |B| - |A \cap B|$，\ 三集合推广公式。
	\item \textbf{逻辑用语}：充分条件与必要条件（$p \Rightarrow q$）、全称量词与存在量词及其否定、命题的四种形式。
\end{itemize}
\end{knowledgeframe}

%% ---- 第一节：集合的基本运算 ---- %%
\section{集合的基本运算}

\startexercise

\subsection*{A组}

\begin{choice}[topic={量词命题的否定}]
已知命题 $p: \exists x \in \mathbb{N}, x^2 \le 0$，则 $\neg p$ 为
\begin{tasks}(2)
	\task $\exists x \in \mathbb{N}, x^2 \le 0$
	\task $\exists x \in \mathbb{N}, x^2 > 0$
	\task $\forall x \in \mathbb{N}, x^2 > 0$
	\task $\forall x \in \mathbb{N}, x^2 \ge 0$
\end{tasks}
\end{choice}
\begin{choice-sol}
\textbf{答案 C}

本题考查全称量词命题与存在量词命题的否定关系.
否定一个量词命题，遵循"改量词，否结论"的原则.
原命题 $p$ 为一个存在量词命题，其逻辑形式为 $\exists x \in M, P(x)$.

其否定的一般形式为：
\[ \neg (\exists x \in M, P(x)) \iff \forall x \in M, \neg P(x) \]
在本题中，量词 $\exists$ (存在) 的否定是 $\forall$ (任意).
原结论 $P(x)$ 为 $x^2 \le 0$，其否定 $\neg P(x)$ 为 $x^2 > 0$.

故 $\neg p$ 为 $\forall x \in \mathbb{N}, x^2 > 0$。

\begin{note}
选项 D 中，$x^2 \ge 0$ 在自然数范围内是恒成立的，
但它并非 $x^2 \le 0$ 的严格否定.
一个结论的否定是其逻辑补集，而非简单的反义关系.
对于实数 $y$，$y \le 0$ 的否定是 $y>0$。
\end{note}
\end{choice-sol}

\begin{choice}[source={2023 秋 · 高一 · 重庆永川区 · 月考校考}, topic={集合的补集与交集}]
设集合 $A=\{x\mid x<2 \text{或} x \ge 4\}$, $B=\{x\mid x<a\}$，若 $(\complement_U A) \cap B \neq \emptyset$，则 $a$ 的取值范围是
\begin{tasks}(2)
	\task $a < 2$
	\task $a > 2$
	\task $a \le 4$
	\task $a \ge 4$
\end{tasks}
\end{choice}
\begin{choice-sol}
\textbf{答案 B}

本题的核心思想是将抽象的集合运算转化为直观的数轴位置关系.
首先，根据集合 $A$ 的定义，确定其在全集 $U=\mathbb{R}$ 下的补集。
\[ A = \{x\mid x<2 \text{ 或 } x \ge 4\} = (-\infty, 2) \cup [4, +\infty) \]
其补集为数轴上未被覆盖的部分：
\[ \complement_U A = \{x\mid 2 \le x < 4\} = [2, 4) \]
集合 $B$ 的定义为：
\[ B = \{x\mid x<a\} = (-\infty, a) \]
题目条件为 $(\complement_U A) \cap B \neq \emptyset$，
即区间 $[2, 4)$ 与 $(-\infty, a)$ 的交集不为空.
为满足此条件，区间 $(-\infty, a)$ 的右端点 $a$
必须大于区间 $[2,4)$ 的左端点 $2$.
因此，可得 $a > 2$。
\end{choice-sol}

\begin{choice}[topic={集合的交集与子集个数}]
已知集合 $A=\{(x,y) \mid x,y \in \mathbb{Z}, \text{且 } xy=4\}$, $B=\{(x,y) \mid x \le y\}$，则 $A \cap B$ 的子集的个数为
\begin{tasks}(4)
	\task 3
	\task 4
	\task 8
	\task 16
\end{tasks}
\end{choice}
\begin{choice-sol}
\textbf{答案 D}

解题分为两个步骤：首先确定交集 $A \cap B$ 的所有元素，然后利用子集个数公式求解.

第一步，列举集合 $A$ 的所有元素.
$A$ 的元素是满足 $xy=4$ 的整数坐标点 $(x,y)$。
\[ A = \{(1,4), (4,1), (-1,-4), (-4,-1), (2,2), (-2,-2)\} \]
第二步，从 $A$ 中筛选出满足集合 $B$ 条件 ($x \le y$) 的元素.
\[ A \cap B = \{(1,4), (-4,-1), (2,2), (-2,-2)\} \]
集合 $A \cap B$ 中含有 4 个元素.

一个含有 $n$ 个元素的有限集合，其子集的个数为 $2^n$.
本题中 $n=4$，故子集个数为 $2^4 = 16$。
\end{choice-sol}


\subsection*{B组}

\begin{choice}[source={2023 秋 · 高一 · 山东烟台 · 月考校考}, topic={充分必要条件}]
下列说法错误的是
\begin{tasks}(1)
	\task ``$A \cap B = B$'' 是 ``$B = \emptyset$'' 的必要不充分条件
	\task ``$x=3$'' 的一个充分不必要条件是 ``$x^2-2x-3=0$''
	\task ``$\lvert x\rvert=1$'' 是 ``$x=1$'' 的必要不充分条件
	\task ``$m$ 是实数'' 的一个充分不必要条件是 ``$m$ 是有理数''
\end{tasks}
\end{choice}
\begin{choice-sol}
\textbf{答案 D}

本题考查充分条件与必要条件的判断.
核心方法是将命题间的逻辑关系转化为集合间的包含关系：
若 $p \Rightarrow q$，则 $p$ 对应的集合是 $q$ 对应集合的子集.

A. ``$A \cap B = B$'' 等价于 $B \subseteq A$.
``$B=\emptyset$'' 能够推出 $B \subseteq A$，但反之不成立.
因此，``$B \subseteq A$'' 是 ``$B=\emptyset$'' 的必要不充分条件.该说法正确.

B. ``$x=3$'' 对应集合 $\{3\}$.
``$x^2-2x-3=0$'' 对应解集 $\{3, -1\}$.
因为 $\{3\} \subsetneq \{3, -1\}$，
所以前者是后者的充分不必要条件.该说法正确.

C. ``$\lvert x\rvert=1$'' 对应集合 $\{1, -1\}$.
``$x=1$'' 对应集合 $\{1\}$.
因为 $\{1\} \subsetneq \{1, -1\}$，
所以前者是后者的必要不充分条件.该说法正确.

D. 命题 $p$: ``$m$ 是有理数''，对应集合 $\mathbb{Q}$.
命题 $q$: ``$m$ 是实数''，对应集合 $\mathbb{R}$.
因为 $\mathbb{Q} \subsetneq \mathbb{R}$，
所以 $p$ 是 $q$ 的充分不必要条件.
原说法颠倒了主次，故错误.
\end{choice-sol}

\begin{choice}[topic={全称命题的否定}]
若命题 ``$\forall x \in \mathbb{R}, 1-x^2 \le m$'' 是假命题，则实数 $m$ 的取值范围是
\begin{tasks}(2)
	\task $(-\infty, 1)$
	\task $(-\infty, 1]$
	\task $(1, +\infty)$
	\task $[1, +\infty)$
\end{tasks}
\end{choice}
\begin{choice-sol}
\textbf{答案 A}

一个全称量词命题为假，其逻辑等价于它的否定形式（一个存在量词命题）为真.
设原命题为 $P: \forall x \in \mathbb{R}, 1-x^2 \le m$。

由于命题 $P$ 为假，其否定命题 $\neg P$ 必为真.
\[ \neg P: \exists x \in \mathbb{R}, \neg(1-x^2 \le m) \implies \exists x \in \mathbb{R}, 1-x^2 > m \]
该存在量词命题为真，意味着 $m$ 必须小于函数 $f(x) = 1-x^2$ 的最大值.
函数 $f(x)=1-x^2$ 是一个开口向下的二次函数，其顶点为最大值点.
\[ f(x)_{\max} = f(0) = 1 \]
因此，实数 $m$ 的取值范围必须是 $m < 1$，
即 $m \in (-\infty, 1)$。
\end{choice-sol}

\begin{choice}[topic={充分不必要条件的应用}]
若 ``$x=2$'' 是 ``$m^2x^2 - (m+3)x + 4 = 0$'' 的充分不必要条件，则实数 $m$ 的值为
\begin{tasks}(2)
	\task 1
	\task $-\frac{1}{2}$
	\task $-\frac{1}{2}$ 或 1
	\task -1 或 $-\frac{1}{2}$
\end{tasks}
\end{choice}
\begin{choice-sol}
\textbf{答案 B}

命题``$p$ 是 $q$ 的充分不必要条件''包含两层含义，必须同时满足：
\begin{itemize}
	\item \textbf{充分性}: $p \Rightarrow q$。
	\item \textbf{不必要性}: $q \not\Rightarrow p$。
\end{itemize}
设 $p: x=2$，$q: m^2x^2 - (m+3)x + 4 = 0$。

\textbf{检验充分性}:
将 $x=2$ 代入方程 $q$ 中，方程必须成立.
\[ 4m^2 - 2(m+3) + 4 = 0 \implies 2m^2 - m - 1 = 0 \]
因式分解得 $(2m+1)(m-1) = 0$，
解得 $m=1$ 或 $m=-\frac{1}{2}$。

\textbf{检验不必要性}:
方程 $q$ 的解集 $S$ 必须真包含 $\{2\}$，即 $S \neq \{2\}$。
\begin{itemize}
	\item 当 $m=1$ 时，方程为 $x^2 - 4x + 4 = 0$，
	解集 $S=\{2\}$.
	此时 $q \Rightarrow p$ 成立，不满足``不必要''条件，故舍去.
	\item 当 $m=-\frac{1}{2}$ 时，
	方程为 $\frac{1}{4}x^2 - \frac{5}{2}x + 4 = 0$，
	即 $x^2 - 10x + 16 = 0$.
	解集 $S=\{2, 8\}$.
	此时 $S \neq \{2\}$，满足``不必要''条件。
\end{itemize}
综上，唯一满足条件的实数 $m$ 的值为 $-\frac{1}{2}$。
\end{choice-sol}

\begin{choice}[topic={集合的补集与互异性}]
设全集 $U=\{1,2,3,4,5\}$，集合 $A=\{1,a,b\}, B=\{4, a-b\}$.若 $\complement_U A = B$，则 $a,b$ 的值分别为
\begin{tasks}(2)
	\task 3,2
	\task 4,3
	\task 3,2 或 5,3
	\task 5,2 或 5,3
\end{tasks}
\end{choice}
\begin{choice-sol}
\textbf{答案 D}

条件 $\complement_U A = B$ 是解题的关键，
它等价于以下两个条件同时成立：
\[ A \cup B = U \quad \text{且} \quad A \cap B = \emptyset \]
根据题意：
$U=\{1,2,3,4,5\}$
$A=\{1,a,b\}$
$B=\{4, a-b\}$

由 $A \cup B = U$，可得：
\[ \{1,a,b\} \cup \{4, a-b\} = \{1, a, b, 4, a-b\} = \{1,2,3,4,5\} \]
比较两集合，可知无序集合 $\{a, b, a-b\}$ 必须与 $\{2,3,5\}$ 相等.
同时，集合 A 和 B 的元素必须满足互异性.

我们通过分类讨论来寻找 $a,b$ 的取值：
\begin{itemize}
	\item \textbf{情况一：} 设 $a=5, b=2$.
	则 $a-b = 3$.
	此时 $\{a,b,a-b\} = \{5,2,3\}$ 与 $\{2,3,5\}$ 匹配.
	检验：$A=\{1,5,2\}, B=\{4,3\}$.
	此时 $\complement_U A = \{3,4\}=B$，符合题意.
	\item \textbf{情况二：} 设 $a=5, b=3$.
	则 $a-b = 2$.
	此时 $\{a,b,a-b\} = \{5,3,2\}$ 与 $\{2,3,5\}$ 匹配.
	检验：$A=\{1,5,3\}, B=\{4,2\}$.
	此时 $\complement_U A = \{2,4\}=B$，符合题意。
\end{itemize}
其他组合（如 $a=3,b=2$ 时 $a-b=1$,与$A$中元素重复）均不满足条件.
因此，$a,b$ 的值为 5,2 或 5,3。
\end{choice-sol}


\subsection*{C组}

\begin{choice}[source={2024 秋 · 高一 · 福建龙岩 · 开学考试校考}, topic={新定义集合与元素互异性}]
定义集合 $A \otimes B = \{x \mid x=\sqrt{a^2+b^2}, a \in A, b \in B\}$，若 $A=\{n, -1\}, B=\{\sqrt{2}, 1\}$，且集合 $A \otimes B$ 中有 3 个元素，则由实数 $n$ 的所有取值组成的集合的非空真子集的个数为
\begin{tasks}(4)
	\task 2
	\task 6
	\task 14
	\task 15
\end{tasks}
\end{choice}
\begin{choice-sol}
\textbf{答案 B}

本题综合考查了新定义集合的理解、集合元素的互异性以及子集个数的计算.

\textbf{1. 构造集合 $A \otimes B$}
根据定义，$A=\{n, -1\}, B=\{\sqrt{2}, 1\}$。
\[ A \otimes B = \{\sqrt{n^2 + (\sqrt{2})^2}, \sqrt{n^2 + 1^2}, \sqrt{(-1)^2 + (\sqrt{2})^2}, \sqrt{(-1)^2 + 1^2}\} \]
化简得：
\[ A \otimes B = \{\sqrt{n^2+2}, \sqrt{n^2+1}, \sqrt{3}, \sqrt{2}\} \]

\textbf{2. 应用元素个数条件}
已知 $\lvert A \otimes B\rvert = 3$.
由于上式给出了 4 个表达式，根据集合元素的互异性，
其中必有两个表达式的值相等.
注意到 $\sqrt{n^2+2} > \sqrt{n^2+1}$ 且 $\sqrt{3} > \sqrt{2}$，
因此相等的只可能是一个含 $n$ 的项与一个常数项.

分类讨论：
\begin{itemize}
	\item 若 $\sqrt{n^2+1} = \sqrt{2} \implies n^2=1 \implies n=\pm 1$.
	此时 $\sqrt{n^2+2}=\sqrt{3}$，
	集合为 $\{\sqrt{2}, \sqrt{3}\}$，
	仅有 2 个元素，不合题意.
	\item 若 $\sqrt{n^2+1} = \sqrt{3} \implies n^2=2 \implies n=\pm \sqrt{2}$.
	此时 $\sqrt{n^2+2}=\sqrt{4}=2$，
	集合为 $\{2, \sqrt{3}, \sqrt{2}\}$，
	有 3 个元素，符合题意.
	\item 若 $\sqrt{n^2+2} = \sqrt{2} \implies n^2=0 \implies n=0$.
	此时 $\sqrt{n^2+1}=1$，
	集合为 $\{1, \sqrt{3}, \sqrt{2}\}$，
	有 3 个元素，符合题意.
	\item 若 $\sqrt{n^2+2} = \sqrt{3} \implies n^2=1$，
	情况同第一类，不符。
\end{itemize}

\textbf{3. 计算子集个数}
由上可知，实数 $n$ 的所有取值组成的集合为
$S_n = \{-\sqrt{2}, 0, \sqrt{2}\}$。
该集合有 $k=3$ 个元素.
非空真子集的个数为 $2^k - 2 = 2^3 - 2 = 6$。
\end{choice-sol}


\stopexercise

%% ---- 第二节：德摩根定律与容斥原理 ---- %%
\section{德摩根定律与容斥原理}

\startexercise

\subsection*{A组}

\begin{prob}[topic={集合的并集与补集运算}, score=5]
\textbf{(填空题)} 已知全集 $U=\{x \in \mathbb{Z} \mid -3 \le x \le 3\}$, 集合 $A=\{-2, 1, 2\}$, $B=\{-1, 0, 1\}$. 求 $\complement_U(A \cup B) = $ \rule{3cm}{0.5pt}.
\end{prob}
\begin{prob-sol}
\textbf{答案} $\{-3, 3\}$

本题考查有限集合的补集运算.

\textbf{方法一：直接法}
首先确定全集 $U$ 和并集 $A \cup B$。
\[ U = \{-3, -2, -1, 0, 1, 2, 3\} \]
\[ A \cup B = \{-2, 1, 2\} \cup \{-1, 0, 1\} = \{-2, -1, 0, 1, 2\} \]
$\complement_U(A \cup B)$ 即为在 $U$ 中但不在 $A \cup B$ 中的元素所构成的集合.
\[ \complement_U(A \cup B) = \{-3, 3\} \]

\textbf{方法二：德摩根定律}
根据德摩根定律，$\complement_U(A \cup B) = (\complement_U A) \cap (\complement_U B)$。
\[ \complement_U A = \{-3, -1, 0, 3\} \]
\[ \complement_U B = \{-3, -2, 2, 3\} \]
两补集的交集为：
\[ (\complement_U A) \cap (\complement_U B) = \{-3, 3\} \]
\end{prob-sol}

\begin{prob}[topic={容斥原理初步应用}, score=10]
\textbf{(解答题)} 某班有50名学生，参加数学兴趣小组的有25人，参加物理兴趣小组的有22人，两个小组都参加的有10人.问：
\begin{enumerate}
	\item 至少参加一个兴趣小组的学生有多少人？
	\item 两个小组都没有参加的学生有多少人？
\end{enumerate}
\end{prob}
\begin{prob-sol}
本题是容斥原理的直接应用.
设参加数学兴趣小组的学生集合为 $M$，
参加物理兴趣小组的学生集合为 $P$。
根据题意，已知信息为：
\[ \lvert M\rvert = 25, \quad \lvert P\rvert = 22, \quad \lvert M \cap P\rvert = 10 \]
(1) \textbf{至少参加一个兴趣小组的人数}，即求 $\lvert M \cup P\rvert$。
根据容斥原理公式：
\[ \lvert M \cup P\rvert = \lvert M\rvert + \lvert P\rvert - \lvert M \cap P\rvert = 25 + 22 - 10 = 37 \text{ (人)} \]

(2) \textbf{两个小组都没有参加的人数}.
设全班学生为全集 $U$，则所求人数为 $\lvert\complement_U(M \cup P)\rvert$。
\[ \lvert\complement_U(M \cup P)\rvert = \lvert U\rvert - \lvert M \cup P\rvert = 50 - 37 = 13 \text{ (人)} \]
\end{prob-sol}


\subsection*{B组}

\begin{prob}[topic={德摩根定律与一元二次不等式}, score=12]
\textbf{(解答题)} 设全集 $U=\mathbb{R}$, 集合 $A=\{x\mid x \le -1 \text{ 或 } x \ge 4\}$, $B=\{x\mid x^2-3x-4<0\}$.
\begin{enumerate}
	\item 求集合 $B$；
	\item 利用德摩根定律，求 $\complement_U(A \cap B)$.
\end{enumerate}
\end{prob}
\begin{prob-sol}
本题将德摩根定律与一元二次不等式的求解相结合.

(1) \textbf{求集合 B}
解不等式 $x^2-3x-4<0$。
因式分解得 $(x-4)(x+1)<0$，解得 $-1 < x < 4$。
\[ B=\{x \mid -1 < x < 4\} = (-1, 4) \]

(2) \textbf{求 $\complement_U(A \cap B)$}
直接计算 $A \cap B$ 再求补集较为繁琐，应用德摩根定律可以简化运算。
\[ \complement_U(A \cap B) = (\complement_U A) \cup (\complement_U B) \]
分别计算两个补集：
\[ \complement_U A = \{x \mid \neg(x \le -1 \text{ 或 } x \ge 4)\} = \{x \mid -1 < x < 4\} = (-1, 4) \]
\[ \complement_U B = \{x \mid \neg(-1 < x < 4)\} = \{x \mid x \le -1 \text{ 或 } x \ge 4\} = (-\infty, -1] \cup [4, \infty) \]
求它们的并集：
\[ (\complement_U A) \cup (\complement_U B) = (-1, 4) \cup \left( (-\infty, -1] \cup [4, \infty) \right) = \mathbb{R} \]
因此，$\complement_U(A \cap B) = \mathbb{R}$。
\end{prob-sol}

\begin{prob}[topic={容斥原理的逆向应用}, score=12]
\textbf{(解答题)} 某次考试中，50名学生参加了数学和物理两科.已知数学成绩及格的有40人，物理成绩及格的有31人，两科都不及格的有4人.问这次考试中：
\begin{enumerate}
	\item 两科都及格的学生有多少人？
	\item 恰好只有一科及格的学生有多少人？
\end{enumerate}
\end{prob}
\begin{prob-sol}
本题是容斥原理的逆向和综合应用.
设数学及格的集合为 $M$，物理及格的集合为 $P$。
已知信息：
\[ \lvert U\rvert=50, \quad \lvert M\rvert=40, \quad \lvert P\rvert=31 \]
``两科都不及格''的人数为 4，即 $\lvert\complement_U(M \cup P)\rvert=4$。

(1) \textbf{求两科都及格的人数}，即求 $\lvert M \cap P\rvert$。
首先，由补集信息计算至少有一科及格的人数 $\lvert M \cup P\rvert$。
\[ \lvert M \cup P\rvert = \lvert U\rvert - \lvert\complement_U(M \cup P)\rvert = 50 - 4 = 46 \]
再根据容斥原理 $\lvert M \cup P\rvert = \lvert M\rvert + \lvert P\rvert - \lvert M \cap P\rvert$，可得：
\[ 46 = 40 + 31 - \lvert M \cap P\rvert \]
\[ \lvert M \cap P\rvert = 71 - 46 = 25 \text{ (人)} \]

(2) \textbf{求恰好只有一科及格的人数}.
这部分学生等于``至少一科及格''的人数减去``两科都及格''的人数。
\[ \lvert(M \cup P) \setminus (M \cap P)\rvert = \lvert M \cup P\rvert - \lvert M \cap P\rvert = 46 - 25 = 21 \text{ (人)} \]
\end{prob-sol}


\subsection*{C组}

\begin{prob}[topic={三集合容斥原理}, score=12]
\textbf{(解答题)} 某新闻机构就A, B, C三个热点话题对100人进行调查，结果显示：关注A的有45人，B的有55人，C的有60人；同时关注A和B的有25人，A和C的有20人，B和C的有30人；三个话题都关注的有10人.问：
\begin{enumerate}
	\item 至少关注一个话题的人数是多少？
	\item 恰好只关注一个话题的人数是多少？
\end{enumerate}
\end{prob}
\begin{prob-sol}
本题是三集合容斥原理的经典应用.
设关注 A, B, C 话题的人构成的集合分别为 $A, B, C$。
已知信息：
\[ \lvert A\rvert=45, \quad \lvert B\rvert=55, \quad \lvert C\rvert=60 \]
\[ \lvert A \cap B\rvert=25, \quad \lvert A \cap C\rvert=20, \quad \lvert B \cap C\rvert=30 \]
\[ \lvert A \cap B \cap C\rvert=10 \]

(1) \textbf{求至少关注一个话题的人数}，即求 $\lvert A \cup B \cup C\rvert$。
应用三集合容斥原理公式：
\begin{align*}
	\lvert A \cup B \cup C\rvert &= (\lvert A\rvert+\lvert B\rvert+\lvert C\rvert) - (\lvert A \cap B\rvert+\lvert A \cap C\rvert+\lvert B \cap C\rvert) + \lvert A \cap B \cap C\rvert \\
	&= (45+55+60) - (25+20+30) + 10 \\
	&= 160 - 75 + 10 \\
	&= 95 \text{ (人)}
\end{align*}

(2) \textbf{求恰好只关注一个话题的人数}.
这部分人数等于只关注 A，只关注 B，只关注 C 的人数之和.

只关注A的人数：
\[ \lvert A \setminus (B \cup C)\rvert = \lvert A\rvert - \lvert A \cap B\rvert - \lvert A \cap C\rvert + \lvert A \cap B \cap C\rvert = 45 - 25 - 20 + 10 = 10 \text{ (人)} \]
只关注B的人数：
\[ \lvert B \setminus (A \cup C)\rvert = \lvert B\rvert - \lvert A \cap B\rvert - \lvert B \cap C\rvert + \lvert A \cap B \cap C\rvert = 55 - 25 - 30 + 10 = 10 \text{ (人)} \]
只关注C的人数：
\[ \lvert C \setminus (A \cup B)\rvert = \lvert C\rvert - \lvert A \cap C\rvert - \lvert B \cap C\rvert + \lvert A \cap B \cap C\rvert = 60 - 20 - 30 + 10 = 20 \text{ (人)} \]

总人数为 $10 + 10 + 20 = 40$ (人).
\end{prob-sol}


\stopexercise
